\pdfmapfile{+symbols.map}
\pdfmapfile{+cm.map}
\pdfmapfile{+cmextra.map}
\documentclass[11pt]{article}
\usepackage[a4paper,margin=25mm]{geometry}
\usepackage{amsmath,amssymb,amsthm}
\usepackage[colorlinks=true,urlcolor=blue]{hyperref}
\newtheorem{theorem}{Theorem}
\newtheorem{lemma}{Lemma}
\title{Records of Exchangeable Permutations\\\large A complete finite counting proof of TLMC 00000009361}
\author{Galaxy-0}
\date{16 October 2026}
\begin{document}
\maketitle
\begin{abstract}
The upper-record indicators of a positionally exchangeable random permutation
are independent Bernoulli variables with parameters $1,1/2,\ldots,1/n$.
We prove the exact joint pattern count using a bijection of actual
permutations. A Lean 4.31 formalization verifies the permutation sample space,
its duplicate-free enumeration, record semantics and the counting identity.
\end{abstract}
\section{Statement and conventions}
The official conjecture describes ``the records of exchangeable permutations''
and says their distribution is independent Bernoulli. We give its standard
precise meaning. Let $\mathcal P_n$ be all permutations of $\{0,\ldots,n-1\}$,
written chronologically as $p=(p_1,\ldots,p_n)$. Define the upper-record
indicator
\[
 R_k(p)=\mathbf 1\{p_k>p_j\text{ for every }j<k\},\qquad 1\leq k\leq n.
\]
The empty comparison makes $R_1=1$. Exchangeability here means invariance
of the law under arbitrary permutations of the \emph{positions}. This is
not merely conjugation invariance when a permutation is viewed as a group
element. We assume distinct entries; ties would change the claim.
\begin{theorem}[Exact joint count]
For every $n\geq0$ and $b=(b_1,\ldots,b_n)\in\{0,1\}^n$,
\begin{equation}\label{eq:count}
 \#\{p\in\mathcal P_n:R_k(p)=b_k\text{ for all }k\}
 =\prod_{k=1}^n \begin{cases}1,&b_k=1,\\ k-1,&b_k=0.\end{cases}
\end{equation}
The empty product is $1$. Thus $\#\mathcal P_n=n!$, and for a positionally
exchangeable permutation the indicators $R_1,\ldots,R_n$ are independent,
with $\Pr(R_k=1)=1/k$.
\end{theorem}
\section{The permutation bijection}
For $0\leq r\leq n$, define the increasing map
\[
 s_r(x)=\begin{cases}x,&x<r,\\x+1,&x\geq r.\end{cases}
\]
It is a bijection from $\{0,\ldots,n-1\}$ to
$\{0,\ldots,n\}\setminus\{r\}$. Its inverse on that set is
\[
 d_r(y)=\begin{cases}y,&y<r,\\y-1,&y>r.\end{cases}
\]
The following is a bijection, with no hidden multiplicities:
\[
 \mathcal P_n\times\{0,\ldots,n\}\longrightarrow\mathcal P_{n+1},\qquad
 (q,r)\longmapsto (s_r(q_1),\ldots,s_r(q_n),r).
\]
Indeed the output contains every label exactly once. Conversely, remove
the last entry $r$ of any permutation in $\mathcal P_{n+1}$ and apply
$d_r$ to every remaining entry. This recovers a unique $q\in\mathcal P_n$.
The identities $d_r(s_r(x))=x$ and $s_r(d_r(y))=y$ for $y\ne r$
prove both inverse identities. Starting with the unique empty permutation,
this also proves $\#\mathcal P_{n+1}=(n+1)\#\mathcal P_n$ and hence $n!$.
\section{Records and the joint count}
Strict monotonicity gives
$s_r(x)<s_r(y)$ if and only if $x<y$. Therefore all the first $n$ record
indicators are preserved by the bijection. The last entry $r$ is a record
if and only if $r=n$: when $r=n$, every previous label is smaller; when
$r<n$, the previous labels include $n$, which is larger.

Fix a record pattern of length $n$. Each realizing permutation consequently
has exactly one extension whose last entry is a record, and exactly $n$
extensions whose last entry is not a record. These extensions are disjoint
and exhaust all permutations with the prescribed extended pattern, by the
bijection above. If $C_n(b)$ is the left-hand side of \eqref{eq:count}, then
\[
 C_{n+1}(b,1)=C_n(b),\qquad C_{n+1}(b,0)=nC_n(b),\qquad C_0(())=1.
\]
Induction proves \eqref{eq:count}. In particular a pattern beginning with
$0$ has count zero, as required; no separate positive-size assumption is used.
\section{Exchangeability and probability}
The position-permutation action on $\mathcal P_n$ is transitive: one can
reorder any permutation into any other. Invariance under this action
therefore assigns the same probability to all $n!$ elements. Since their
probabilities sum to one, each has probability $1/n!$. Dividing
\eqref{eq:count} by $n!=\prod_{k=1}^n k$ gives
\[
 \Pr(R_1=b_1,\ldots,R_n=b_n)
 =\prod_{k=1}^n
 \begin{cases}1/k,&b_k=1,\\1-1/k,&b_k=0.\end{cases}
\]
This equality for every Boolean pattern is precisely mutual independence
and the claimed Bernoulli laws. The \emph{number} of records is the sum
$N_n=\sum_{k=1}^nR_k$, not a family of independent copies of itself. Its
probability generating function is
\[
 \mathbb E[z^{N_n}]=\prod_{k=1}^n\left(1-\frac1k+\frac zk\right).
\]
At $n=0$, the law is concentrated on $N_0=0$. For any other fixed ordered
set of $n$ distinct labels, increasing relabeling preserves all comparisons
and yields the same result.
\clearpage
\section{Formalization and trust boundary}
The source uses only Lean 4.31.0 and its bundled Std library.
\texttt{Permutations.lean} defines actual lists of natural-number labels.
The theorem \texttt{perms\_complete} identifies membership of its enumeration
with \texttt{List.Perm (List.range n)}; \texttt{perms\_nodup} rules out
repeated sample points and \texttt{perms\_length} proves its factorial size.
The intermediate validity predicate requires both no repeated entries
and the exact membership condition $x\in p\iff x<n$.

For convenient recursion the enumeration stores a permutation in reverse
chronological order. This is only a storage convention:
\texttt{RecordSemantics.lean} separately defines a left-to-right scan which
compares each entry with every earlier entry.
\texttt{recs\_chronological} proves agreement after reversal.
\texttt{recs\_extend} proves preservation of all earlier records and the
last-record criterion $r=n$. \texttt{RecordsCount.lean} proves \texttt{chronologicalPatternCount\_exact} for the complete duplicate-free chronological enumeration. Its \texttt{independent\_bernoulli\_joint\_law} also proves that the exact rational uniform probability is the Bernoulli product. These are statements about actual permutations, rather than an unconnected sequence of rank choices.

The finite counting statements, normalized rational product identity and record/bijection correspondence are kernel-checked. The implication from positional exchangeability to the uniform measure and the interpretation of the joint product as independence are the elementary arguments in the preceding section; no general probability-library formalization is claimed.
No admitted proofs, custom conclusion axioms, native-decision shortcuts or
kernel bypasses are used. The audit file records the theorem axiom footprints.
A separate Python test directly enumerates chronological permutations for
$n=0,\ldots,8$ and checks every Boolean pattern. These tests are supplemental,
not assumptions of the Lean proof.
\section*{Source}
\href{https://github.com/The-Last-Math-Competition/The-Last-Math-Competition/blob/main/conjectures/00000009361.md}{Official TLMC conjecture 00000009361}.
\par\smallskip
Source blob SHA: \texttt{3c1e3568c9aa81b24440aa69a23149d953e1b4d4}.
The English and Chinese statement is preserved in the accompanying source file.
\end{document}
